\uuid{190d}
\exo7id{4168}
\titre{exo7 4168}
\auteur{quercia}
\organisation{exo7}
\datecreate{2010-03-11}
\isIndication{false}
\isCorrection{false}
\chapitre{Fonction de plusieurs variables}
\sousChapitre{Dérivée partielle}
\cpge{mpsi-mp2i.fonctions-deux-variables.b, pcsi.fonctions-deux-variables.b, ptsi.fonctions-deux-variables.b, bcpst1.fonctions-deux-variables.c, mp-mpi.calcul-differentiel-optimisation.b, pc.fonctions-plusieurs-variables.c, psi.fonctions-plusieurs-variables.c}
\module{Analyse}
\niveau{L2}
\difficulte{}

\contenu{
\texte{
Soit $f:\R \to \R$ de classe $\mathcal{C}^1$.
}
\begin{enumerate}
    \item \question{Montrer que, sous une condition à préciser, l'équation $y-zx=f(z)$ définit
    localement $z$ fonction implicite de $x$ et $y$.}
    \item \question{Montrer que l'on a alors : $\frac{\partial z}{\partial x} + z\frac{\partial z}{\partial y} = 0$.}
\end{enumerate}
}
